% from ega0-1-fr.tex
\oldpage[0\textsubscript{I}]{55}

Si $M$, $N$ sont deux $A$-modules, $M'$ (resp. $N'$) un sous-module de $M$
(resp. $N$), nous noterons $\operatorname{Im}(M'\otimes_A N')$ le sous-module
de $M\otimes_A N$, image de l'application canonique
$M'\otimes_A N'\to M\otimes_A N$.
\end{env}

\subsection{Modules plats}
\label{subsection:0.6.1-fr}

\begin{env}[6.1.1]
\label{0.6.1.1-fr}
Soit $M$ un $A$-module. Les propriétés suivantes sont équivalentes~:
\begin{enumerate}
  \item[$a)$] Le foncteur $M\otimes_A N$ en $N$ est exact dans la catégorie
    des $A$-modules~;
  \item[$b)$] $\operatorname{Tor}_i^A(M,N)=0$ pour tout $i>0$ et tout
    $A$-module $N$~;
  \item[$c)$] $\operatorname{Tor}_1^A(M,N)=0$ pour tout $A$-module $N$.
\end{enumerate}

Lorsque $M$ vérifie ces conditions, on dit que $M$ est un
\emph{$A$-module plat}. Il est clair que tout $A$-module libre est plat.

Pour que $M$ soit un $A$-module plat, il suffit que pour tout idéal
$\mathfrak{J}$ de $A$, \emph{de type fini}, l'application canonique
$M\otimes_A\mathfrak{J}\to M\otimes_A A=M$ soit \emph{injective}.
\end{env}

\begin{env}[6.1.2]
\EGAReaderTarget{canon-ref.0.6.1.2-fr}{0.6.1.2-fr}
Toute limite inductive de $A$-modules plats est un $A$-module plat. Pour
qu'une somme directe $\bigoplus_{\lambda\in L}M_\lambda$ de $A$-modules soit
un $A$-module plat, il faut et il suffit que chacun des $A$-modules
$M_\lambda$ soit plat. En particulier, tout $A$-module projectif est plat.

Soit
\[
  0\longrightarrow M'\longrightarrow M\longrightarrow M''\longrightarrow 0
\]
une suite exacte de $A$-modules, telle que $M''$ soit \emph{plat}. Alors, pour
tout $A$-module $N$, la suite
\[
  0\longrightarrow M'\otimes N\longrightarrow M\otimes N
  \longrightarrow M''\otimes N\longrightarrow 0
\]
est exacte. En outre, pour que $M$ soit plat, il faut et il suffit alors que
$M'$ le soit (mais il peut se faire que $M$ et $M'$ soient plats sans que
$M''=M/M'$ le soit).
\end{env}

\begin{env}[6.1.3]
\label{0.6.1.3-fr}
Soient $M$ un $A$-module plat, $N$ un $A$-module quelconque~; pour deux
sous-modules $N'$, $N''$ de $N$, on a alors
\[
  \operatorname{Im}(M\otimes(N'+N''))
  =\operatorname{Im}(M\otimes N')+\operatorname{Im}(M\otimes N''),
\]
\[
  \operatorname{Im}(M\otimes(N'\cap N''))
  =\operatorname{Im}(M\otimes N')\cap\operatorname{Im}(M\otimes N'')
\]
(images prises dans $M\otimes N$).
\end{env}

\begin{env}[6.1.4]
\label{0.6.1.4-fr}
Soient $M$, $N$ deux $A$-modules, $M'$ (resp. $N'$) un sous-module de $M$
(resp. $N$), et supposons que l'un des modules $M/M'$, $N/N'$ soit plat.
Alors on a
\[
  \operatorname{Im}(M'\otimes N')
  =\operatorname{Im}(M'\otimes N)\cap\operatorname{Im}(M\otimes N')
\]
(images dans $M\otimes N$). En particulier, si $\mathfrak{J}$ est un idéal
de $A$ et si $M/M'$ est plat, on a
$\mathfrak{J}M'=M'\cap\mathfrak{J}M$.
\end{env}

\subsection{Changement d'anneaux}
\label{subsection:0.6.2-fr}

Lorsqu'un groupe additif $M$ est muni de plusieurs structures de module par
rapport à des anneaux $A$, $B$, \ldots, au lieu de dire que $M$ est plat en
tant que $A$-module, $B$-module, \ldots, nous dirons parfois aussi que $M$
est \emph{$A$-plat}, \emph{$B$-plat}, \ldots

\begin{env}[6.2.1]
\label{0.6.2.1-fr}
Soient $A$ et $B$ deux anneaux, $M$ un $A$-module, $N$ un
$(A,B)$-bimodule. Si $M$ est plat et si $N$ est $B$-plat, alors
$M\otimes_A N$ est $B$-plat. En particulier, si $M$ et $N$ sont deux
$A$-modules plats, $M\otimes_A N$ est un $A$-module plat. Si $B$ est une
$A$-algèbre et si $M$ est
